\chapter{System Overview}

\section{Main Stakeholders}
\defaultauthor

Các bên liên quan gồm du khách, nhóm phát triển, nhà cung cấp dữ liệu/dịch vụ và hội đồng đánh giá học thuật. IMOVE sử dụng OneMap, LTA DataMall, OpenWeather, Gemini, Supabase và OpenStreetMap để hỗ trợ quyết định; hệ thống không trực tiếp điều phối hoạt động của đơn vị vận tải.

\section{Main Actors}
\defaultauthor

\begin{table}[H]
\centering
\begin{tabularx}{\textwidth}{p{3.5cm}X X}
\toprule
\textbf{Actor} & \textbf{Nhu cầu} & \textbf{Hỗ trợ từ IMOVE} \\
\midrule
Du khách lần đầu & Khó chọn thứ tự và mode & Lập lịch đa điểm, tối ưu và hướng dẫn theo leg \\
Người có preference & Muốn cân bằng thời gian, chi phí, walking, transfers & Weighted profile và constraints \\
Người đang đi & Cần phản ứng trước mưa hoặc disruption & Alert và adaptation proposal \\
Người đăng nhập & Muốn lưu trip và preference & Supabase Auth và persistence \\
\bottomrule
\end{tabularx}
\caption{Các actor chính của hệ thống.}
\end{table}

\section{Core Features}
\defaultauthor

\begin{enumerate}
  \item Lập lịch trình du lịch nhiều ngày và nhiều địa điểm.
  \item Tối ưu thứ tự bằng greedy heuristic có time-window constraints.
  \item So sánh và lựa chọn Metro, Bus, Walk, Cycle và Grab.
  \item Chấm điểm thích nghi theo preference, mưa và giờ cao điểm.
  \item Thích nghi với cảnh báo LTA và thời tiết.
  \item Học preference từ feedback lặp lại.
  \item Gemini hỗ trợ NLP, suggestion, warning và chatbot.
\end{enumerate}

\section{Innovation Highlights}
\defaultauthor

\begin{itemize}
  \item \textbf{Hybrid planning:} kết hợp code xác định, external APIs và LLM.
  \item \textbf{Context-aware recommendation:} điều chỉnh scoring theo bối cảnh thực tế.
  \item \textbf{Graceful degradation:} giữ kết quả có ích khi một số mode không khả dụng.
  \item \textbf{Human-in-the-loop adaptation:} chỉ persist proposal khi người dùng chấp nhận.
  \item \textbf{Explainable learning:} quy tắc học preference đơn giản và kiểm thử được.
\end{itemize}

\section{Overall Architecture}
\defaultauthor

\begin{figure}[H]
\centering
\begin{tikzpicture}[node distance=1.2cm and 1.2cm]
  \node[imovenode] (user) {Traveller};
  \node[imovenode,right=of user] (web) {React + Vite\\Frontend};
  \node[imovenode,right=of web] (api) {FastAPI\\Routers};
  \node[imoveagent,above right=0.5cm and 1.1cm of api] (agents) {Planning / Adaptation\\Memory / Chat};
  \node[imoveagent,below right=0.5cm and 1.1cm of api] (services) {Routing / Transit\\Weather / AI / Scoring};
  \node[imovedata,right=of agents] (db) {Supabase};
  \node[imoveexternal,right=of services] (ext) {External APIs};
  \draw[imovearrow] (user) -- (web);
  \draw[imovearrow] (web) -- (api);
  \draw[imovearrow] (api) -- (agents);
  \draw[imovearrow] (api) -- (services);
  \draw[imovearrow] (agents) -- (db);
  \draw[imovearrow] (services) -- (ext);
\end{tikzpicture}
\caption{Kiến trúc tổng quan của IMOVEV2.}
\label{fig:overall-architecture}
\end{figure}
